\chapter{Resonance and Meaning Curvature}\label{ch:resonance}

An idea algebra is a tuple $(\Ideas,\compose,\unit,\rho,\sigma,\nu,\eta,\omega,\deg)$ in the notation of the Foundations chapter. The previous chapters built the carrier $\Ideas$ and equipped it with an
associative composition $\compose$ with two-sided unit $\unit$. From that
skeleton alone little can be said about \emph{how} ideas relate to one
another. The present chapter introduces the second pillar of the formalism:
the \textbf{resonance pairing} $\rho:\Ideas\times\Ideas\to\RR$ and the
\textbf{meaning curvature} $\curv$ that measures its anti-symmetric part.
Resonance is the quantitative scalar against which every later theorem of
Volume~1 is weighed; curvature is the homological invariant that tells us
when two idea algebras are the \emph{same} algebra or only superficially
similar.

The development follows the Lean files
\texttt{IdeaTheory.Theorems2}, \texttt{IdeaTheory.Theorems4}, and
\texttt{IdeaTheory.Theorems6}, restated in plain mathematical prose with
math-prose proofs. We use the abbreviation $\res{a}{b}$ for $\rho(a,b)$ when
the pairing is unambiguous; we write $V(\Ideas):=\RR^{(\Ideas)}$ for the
free $\RR$-vector space on the underlying set $\Ideas$.

The chapter is organized as follows. \cref{sec:why-pairing} motivates the
bilinearity of $\rho$ from three independent considerations.
\cref{sec:free-vs} constructs the free $\RR$-vector space on $\Ideas$ and
proves the bilinear extension theorem with full uniqueness.
\cref{sec:nondeg} treats non-degeneracy and its preservation under
composition. \cref{sec:sym-antisym} establishes the canonical
decomposition $\rho=\rho_++\rho_-$. \cref{sec:curvature} introduces the
meaning curvature as the antisymmetric part and proves its closedness.
\cref{sec:cohomology} establishes the cohomological invariance of
$[\curv]$ under structural equivalence. \cref{sec:aufhebung} verifies
compatibility with the Aufhebung decomposition. \cref{sec:resonance-zoo}
exhibits seven entries from the zoo of historical and scientific
positions about ideas, each faithfully captured by some component of
the formalism.

% ---------------------------------------------------------------
\section{From dyadic affinity to a bilinear form}\label{sec:why-pairing}

The earliest mathematically articulate doctrine of how one idea ``calls
forth'' another is Hume's account of the \emph{association of ideas}: every
idea, on appearing in the mind, exerts a kind of ``gentle force'' upon
related ideas, drawing them into consciousness in proportion to the strength
of their connexion~\cite{Hume1739}. Hume named three principles of
connexion --- resemblance, contiguity, causation --- and asserted that the
strength was, at least in principle, scalar and symmetric.

\begin{interpretation}\label{int:hume-resonance}
Hume's ``gentle force'' is what we will formalize as the resonance pairing:
a function $\rho:\Ideas\times\Ideas\to\RR$ assigning to each ordered pair
$(a,b)$ a real number measuring the degree to which entertaining $a$
\emph{calls forth} $b$. The asymmetry that Hume neglected --- that
``Newton'' calls ``apple'' more readily than ``apple'' calls ``Newton'' ---
will turn out to be the carrier of the curvature $\curv$.
\end{interpretation}

What forces $\rho$ to be \emph{bilinear} when extended to formal sums? Three
distinct lines of reasoning converge on the same conclusion.

\paragraph{(i) Aggregation of evidence.}
If a population of agents, each holding some idea $a_i$ with weight
$\lambda_i$, encounters a probe $b$, the population-level resonance is the
weighted sum $\sum_i \lambda_i\,\rho(a_i,b)$. Replacing the implicit mixture
$\sum_i\lambda_i a_i$ in the first slot of $\rho$ by anything other than the
weighted sum would force the formalism to specify how the agents'
contributions interact \emph{before} the probe $b$ is even known --- a
context-dependence that contradicts the dyadic character of the original
relation~\cite{Sperber1996}.

\paragraph{(ii) Geometric realization.}
Every concrete instantiation of idea-affinity in the cognitive sciences
that has survived empirical scrutiny --- conceptual
spaces~\cite{Gardenfors2000}, holographic reduced
representations~\cite{Plate1995}, distributed word
embeddings~\cite{Mikolov2013} --- realizes ideas as vectors and resonance
as an inner product (or a closely related bilinear form). To accept any of
these as a model of $\Ideas$ is to commit, in advance, to bilinearity of
$\rho$ on the free vector space.

\paragraph{(iii) Internal coherence.}
The composition $\compose$ is the only resource the bare algebra provides
for combining ideas. The resonance of a composite $a\compose b$ against a
probe $c$ admits, by the $\compose$-decomposition theorem of
Chapter~\ref{ch:resonance} (verified as
\texttt{IdeaTheory.Theorems2.rs\_op\_decomp}), a canonical splitting
$\rho(a\compose b,c)=\rho(a,c)+\rho(b,c)+\omega(a,b,c)$ where $\omega$ is
real-valued. \emph{Almost-additivity} of $\rho$ in the first slot, modulo
the residue $\omega$, is a meaningless property unless $\rho$ is
\emph{actually} additive on the free vector space; the residue $\omega$
becomes the 2-cocycle of Chapter~\ref{ch:cocycle}.

These three motivations all push $\rho$ in the same direction: linearity in
each slot. The remaining freedom is whether linearity is over $\RR$ alone
or over a richer ground field. Because the only quantitative axiom we have
imposed --- $\rho(a,a)\ge 0$ with equality only at $a=\unit$ --- is a
\emph{real} statement, $\RR$ is the canonical and most parsimonious choice.

It is worth recording why \emph{any} bilinear extension at all is the right
move, given that on the carrier set $\Ideas$ there is no a priori
addition. The free vector space construction is precisely the universal
device that supplies addition where none was given, in such a way that all
the structure already present on $\Ideas$ (here: the function $\rho_0$ and
the composition $\compose$) extends uniquely. Without this passage to
$V(\Ideas)$, every quantitative statement of the form ``the resonance of a
mixture is the mixture of the resonances'' would require an axiomatic
postulate; with it, no postulate is needed: the universal property does
the work. The book takes this as a foundational stance, because every
applied chapter --- from cultural transmission (Volume~2) to evidence
accumulation (Volume~4) --- routinely needs to evaluate $\rho$ on
mixtures, often without remarking on the fact.

A subtler point is that $V(\Ideas)$ is \emph{not} itself an idea algebra
in the strict sense. The composition $\compose$ extends bilinearly to $V$
(define $[a]\compose[b]:=[a\compose b]$ on the basis and extend), but the
extension is associative only on the image of $\Ideas$; off the image,
associativity is automatic but unitality is restricted to the
distinguished basis vector $[\unit]$. We will use $V$ throughout as the
\emph{ambient} linear space against which the various bilinear and
multilinear data of Idea Theory are most naturally formulated, while
continuing to think of the genuine ideas as the points of $\Ideas$.

% ---------------------------------------------------------------
\section{The free vector space and bilinear extension}\label{sec:free-vs}

We now make the extension precise. Throughout, $\Ideas$ is a set;
$V:=\RR^{(\Ideas)}$ is the $\RR$-vector space of finitely supported
functions $\Ideas\to\RR$, with canonical basis $\{[a]:a\in\Ideas\}$ where
$[a]$ is the indicator of $\{a\}$.

\begin{definition}[Free vector space on $\Ideas$]\label{def:free-vs}
$V(\Ideas)$ is the universal $\RR$-vector space generated by the set
$\Ideas$: there is an injection $\iota:\Ideas\hookrightarrow V$, $a\mapsto
[a]$, such that for every $\RR$-vector space $W$ and every set map
$f:\Ideas\to W$ there is a unique $\RR$-linear map $\widetilde f:V\to W$
with $\widetilde f\circ\iota=f$.
\end{definition}

\begin{definition}[Resonance structure]\label{def:res-struct}
A \emph{resonance structure} on the bare idea algebra $\IdMon$ is a
function $\rho_0:\Ideas\times\Ideas\to\RR$ satisfying:
\begin{enumerate}
  \item \textbf{Boundary at the unit:} $\rho_0(\unit,a)=\rho_0(a,\unit)=0$
        for every $a\in\Ideas$;
  \item \textbf{Self-positivity:} $\rho_0(a,a)\ge 0$ for every $a$, with
        equality iff $a=\unit$.
\end{enumerate}
\end{definition}

\noindent (Self-positivity is verified as
\texttt{IdeaTheory.Theorems2.weight\_eq\_zero\_iff}.)

\begin{theorem}[Bilinear extension]\label{thm:bilinear-extension}
For every resonance structure $\rho_0$ on $\Ideas$, there exists a unique
bilinear form
\[
\rho:V(\Ideas)\times V(\Ideas)\longrightarrow \RR
\]
such that $\rho([a],[b])=\rho_0(a,b)$ for all $a,b\in\Ideas$.
\end{theorem}

\begin{proof}
\emph{Existence.} For each fixed $b\in\Ideas$, the map
$a\mapsto\rho_0(a,b)$ is a function $\Ideas\to\RR$. By the universal
property of the free vector space, it extends uniquely to an $\RR$-linear
functional $\rho_0(\cdot,b):V\to\RR$. The assignment
$b\mapsto\rho_0(\cdot,b)$ is then a function $\Ideas\to V^*$. Applying the
universal property again, this extends uniquely to an $\RR$-linear map
$\widehat\rho:V\to V^*$. Define $\rho(u,v):=(\widehat\rho(v))(u)$. By
construction $\rho$ is linear in $v$; linearity in $u$ follows because
$\widehat\rho(v)\in V^*$ is itself linear.

\emph{Uniqueness.} Any bilinear form $\rho$ on $V$ that agrees with $\rho_0$
on the basis $\{[a]\}$ is determined on a general pair
$u=\sum_i \lambda_i[a_i]$, $v=\sum_j \mu_j[b_j]$ by
$\rho(u,v)=\sum_{i,j}\lambda_i\mu_j\,\rho_0(a_i,b_j)$, which depends only on
$\rho_0$.\qed
\end{proof}

\begin{remark}[Notation]
We continue to write $\rho$ for the extended form on $V\times V$ as well as
for its restriction to $\Ideas\times\Ideas$. The pairing $\res{u}{v}$ is
shorthand for $\rho(u,v)$.
\end{remark}

\begin{proposition}[Functoriality]\label{prop:functorial}
Let $\phi:\IdMon\to(\Ideas',\compose',\unit')$ be a homomorphism of bare
idea algebras. If $\rho_0,\rho_0'$ are resonance structures with
$\rho_0'(\phi(a),\phi(b))=\rho_0(a,b)$ for all $a,b$, then the linear
extension $\widetilde\phi:V(\Ideas)\to V(\Ideas')$ satisfies
$\rho'(\widetilde\phi u,\widetilde\phi v)=\rho(u,v)$.
\end{proposition}
\begin{proof}
Both sides are bilinear in $(u,v)$ and agree on the basis pairs
$([a],[b])$ by hypothesis; uniqueness in
Theorem~\ref{thm:bilinear-extension} forces the equality.\qed
\end{proof}

% ---------------------------------------------------------------
\section{Non-degeneracy and its preservation}\label{sec:nondeg}

\begin{definition}[Non-degeneracy]\label{def:nondeg}
The form $\rho:V\times V\to\RR$ is \emph{left non-degenerate} if
$\rho(u,v)=0$ for all $v\in V$ implies $u=0$, and \emph{right
non-degenerate} if $\rho(u,v)=0$ for all $u\in V$ implies $v=0$. It is
\emph{non-degenerate} if both hold.
\end{definition}

The self-positivity axiom of Definition~\ref{def:res-struct} is already
strong enough to deliver one half of non-degeneracy on the level of basis
elements:

\begin{lemma}[Pointwise non-vanishing]\label{lem:pointwise-nondeg}
For every $a\in\Ideas\setminus\{\unit\}$, $\rho(a,a)>0$. Hence the family
$\{[a]\}_{a\ne\unit}$ contains no isotropic vector.
\end{lemma}

\begin{proof}
By Definition~\ref{def:res-struct}(2), $\rho_0(a,a)\ge 0$ with equality iff
$a=\unit$.\qed
\end{proof}

\begin{theorem}[Non-degeneracy on the free space]\label{thm:nondeg}
Let $L:=\{u\in V:\rho(u,v)=0\text{ for all }v\in V\}$ be the left radical.
Then $[a]\notin L$ for any $a\ne\unit$. If, in addition, $\rho_0$ is
\emph{positive semi-definite} on $V$ (i.e.\ $\rho(u,u)\ge 0$ for all $u\in
V$ and $\rho(u,u)=0\Rightarrow u\in L$), then the symmetric part $\rho_+$
introduced in \cref{sec:sym-antisym} restricts to a positive-definite form
on the quotient $V/L$, and $L=\{0\}$ exactly when this quotient is all of
$V$.
\end{theorem}
\begin{proof}
The first claim is Lemma~\ref{lem:pointwise-nondeg}. For the second, the
symmetric form $\rho_+(u,v):=\tfrac12(\rho(u,v)+\rho(v,u))$ has
$\rho_+(u,u)=\rho(u,u)\ge 0$. The kernel of $\rho_+$ on $V$ is contained in
$L$ by the polarization identity, and the residue is positive-definite by
construction.\qed
\end{proof}

\begin{theorem}[Preservation under composition]\label{thm:nondeg-comp}
The free vector space $V$ carries a natural $\RR$-bilinear extension of
$\compose$, defined on basis elements by $[a]\compose[b]:=[a\compose b]$ and
extended bilinearly. If $\rho$ is non-degenerate and the linear extension
of $\compose$ is left-cancellative on $V$, then for every $c\in\Ideas$ the
form $\rho_c(u,v):=\rho(u\compose [c],v)$ is again non-degenerate.
\end{theorem}

\begin{proof}
Let $u\in V$ satisfy $\rho_c(u,v)=0$ for all $v$. Then the vector
$u\compose[c]\in V$ pairs to zero against every $v$, so $u\compose[c]\in L$.
By non-degeneracy of $\rho$, $u\compose[c]=0$. Left cancellation on $V$
gives $u=0$.\qed
\end{proof}

\begin{remark}
In the cognitive-science models of \cref{sec:sym-antisym} (HRR, conceptual
spaces, embeddings), the linear extension of $\compose$ is realized either
as tensor product or as circular convolution; left cancellation holds in
both cases on a dense subspace, and Theorem~\ref{thm:nondeg-comp} explains
why these models do not collapse to triviality after a single
binding-and-probing operation.
\end{remark}

We pause to make explicit how strong the non-degeneracy hypothesis is, and
when it can be expected to hold. Three regimes are typical.

\paragraph{Finite-rank regime.} If $\Ideas$ is finite of cardinality $n$,
then $V(\Ideas)\cong\RR^n$ and $\rho$ is represented by an $n\times n$
real matrix $M$ with $M_{ab}=\rho(a,b)$. Non-degeneracy of $\rho$ is
$\det M\ne 0$, an open condition. The self-positivity axiom forces the
diagonal entries of $M$ to be strictly positive at every $a\ne\unit$, but
imposes no a priori bound on the off-diagonal entries; nothing prevents
$\det M=0$ in the absence of further hypotheses such as the
positive-semidefiniteness invoked in \cref{thm:nondeg}.

\paragraph{Hilbertian regime.} If $\Ideas$ embeds in a real Hilbert space
$H$ via $\iota:\Ideas\hookrightarrow H$ with $\rho(a,b)=\langle\iota
a,\iota b\rangle$, then $\rho$ is automatically symmetric and positive
semi-definite, and non-degeneracy reduces to the linear independence of
$\{\iota a:a\ne\unit\}$ in $H$. This is the regime of conceptual spaces
and word embeddings. Curvature is identically zero in this regime --- a
limitation we will lift in \cref{sec:sym-antisym}.

\paragraph{Generic regime.} In the absence of any geometric realization,
non-degeneracy of $\rho$ is a substantive structural assumption to be
imposed when needed. The Lean development records non-degeneracy as a
\emph{property} of a resonance structure, not as part of the basic
definition. Theorems requiring it are flagged accordingly, and the
preservation result (\cref{thm:nondeg-comp}) shows that the property is
inherited by composition with any element of $\Ideas$.

In what follows, every theorem is stated in the regime of greatest
generality compatible with its content; non-degeneracy is invoked only
when needed.

% ---------------------------------------------------------------
\section{Symmetric and antisymmetric parts}\label{sec:sym-antisym}

The form $\rho$ need not be symmetric. The decomposition into symmetric and
antisymmetric parts is canonical and exhaustive.

\begin{definition}[Symmetric/antisymmetric parts]\label{def:sym-asym}
For $u,v\in V$ define
\[
\rho_+(u,v):=\tfrac{1}{2}\bigl(\rho(u,v)+\rho(v,u)\bigr),\qquad
\rho_-(u,v):=\tfrac{1}{2}\bigl(\rho(u,v)-\rho(v,u)\bigr).
\]
\end{definition}

\begin{proposition}[Decomposition]\label{prop:sym-decomp}
$\rho=\rho_+ + \rho_-$, with $\rho_+$ symmetric and $\rho_-$ antisymmetric.
The decomposition is unique: if $\rho=\sigma+\alpha$ with $\sigma$
symmetric and $\alpha$ antisymmetric, then $\sigma=\rho_+$ and
$\alpha=\rho_-$.
\end{proposition}

\begin{proof}
Direct expansion gives $\rho_+(u,v)=\rho_+(v,u)$ and
$\rho_-(u,v)=-\rho_-(v,u)$. For uniqueness: from $\sigma+\alpha=\rho$ and
$\sigma(v,u)+\alpha(v,u)=\rho(v,u)$ we obtain
$\sigma(u,v)=\tfrac12(\rho(u,v)+\rho(v,u))=\rho_+(u,v)$ on adding, and
$\alpha(u,v)=\rho_-(u,v)$ on subtracting.\qed
\end{proof}

The two parts have radically different mathematical and interpretive
characters.

\begin{interpretation}\label{int:sym-asym}
$\rho_+$ measures \emph{shared content}: the magnitude of $\rho_+(a,b)$ is
how much $a$ and $b$ have in common, regardless of which is the
``stimulus'' and which the ``response''. By contrast $\rho_-$ measures
\emph{directional bias}: the sign of $\rho_-(a,b)$ tells us whether $a$
calls forth $b$ more readily than $b$ calls forth $a$, and the magnitude
quantifies the asymmetry. The Humean tradition acknowledges only $\rho_+$;
the Sperberian, Bakhtinian, and Lakovian traditions all require $\rho_-$.
\end{interpretation}

The duality of $\rho_+$ and $\rho_-$ has empirical correlates that long
predate any of the formalisms in our zoo. Tversky's classical
psychometric study of similarity judgments showed that subjects
systematically rated ``North Korea is similar to China'' more highly than
``China is similar to North Korea'' --- a robust effect across many
stimulus pairs. Tversky modelled the effect by a feature-contrast model
with deliberately asymmetric weights; in our framework the asymmetry is
not a model parameter but a coordinate of the bilinear form, namely
$\rho_-(\text{N.\ Korea},\text{China})>0$. The Humean reduction
$\rho=\rho_+$ is then not merely incomplete but empirically false on a
wide range of natural similarity judgments. The discipline-spanning
significance of this point cannot be overstated: every formalism that
identifies similarity with an inner product (and there are many) is
committed to $\rho_-=0$, and is therefore committed to predicting that
similarity is symmetric --- a prediction that is repeatedly falsified.

\begin{lemma}[Polarization]\label{lem:polarization}
For $u,v\in V$,
\[
\rho_+(u,v) \;=\; \tfrac14\bigl(\rho(u+v,u+v)-\rho(u-v,u-v)\bigr).
\]
\end{lemma}

\begin{proof}
Expand $\rho(u\pm v,u\pm v)=\rho(u,u)\pm\rho(u,v)\pm\rho(v,u)+\rho(v,v)$,
subtract, and divide by $4$.\qed
\end{proof}

\begin{lemma}[Diagonal identification]\label{lem:diag}
For all $a\in\Ideas$, $\rho_-(a,a)=0$ and $\rho_+(a,a)=\rho(a,a)$.
Consequently the self-resonance $\mathrm{wt}(a):=\rho(a,a)$ is detected
entirely by $\rho_+$.
\end{lemma}

\begin{proof}
Antisymmetry forces $\rho_-(a,a)=-\rho_-(a,a)$, hence $\rho_-(a,a)=0$. The
second follows from $\rho=\rho_++\rho_-$.\qed
\end{proof}

\begin{lemma}[Cauchy--Schwarz for the symmetric part]\label{lem:cs}
Suppose $\rho_+$ is positive semi-definite on $V$. Then for all $u,v\in V$,
\[
\rho_+(u,v)^2 \;\le\; \rho_+(u,u)\,\rho_+(v,v).
\]
Equality holds iff $u$ and $v$ are linearly dependent modulo the radical
of $\rho_+$.
\end{lemma}
\begin{proof}
Standard: the discriminant of the non-negative quadratic
$t\mapsto\rho_+(u+tv,u+tv)$ in $t$ is non-positive.\qed
\end{proof}

This Cauchy--Schwarz inequality is the formal source of every cosine
similarity that appears in the cognitive-science zoo entries below: when
$\rho_+(a,a)=\rho_+(b,b)=1$ (a normalization condition), the
\emph{cosine similarity} $\rho_+(a,b)$ is bounded in $[-1,1]$, and the
familiar interpretation of cosine as ``aligned'' versus ``orthogonal''
versus ``opposite'' is recovered.

The antisymmetric part has no analogous bound: $\rho_-(a,b)$ can in
principle grow without limit relative to the diagonal, and large values
of $|\rho_-(a,b)|$ are not pathological but rather signal strong
directionality. The interplay of bounded $\rho_+$ and unbounded
$\rho_-$ underlies most of the empirical phenomenology of asymmetric
similarity in the psychometric literature.

% ---------------------------------------------------------------
\section{The meaning curvature $\curv$}\label{sec:curvature}

The antisymmetric part $\rho_-$ deserves a name of its own. Following the
analogy with differential geometry --- where the antisymmetric part of a
connection 1-form is the curvature 2-form --- we call it the
\textbf{meaning curvature}.

\begin{definition}[Meaning curvature]\label{def:curvature}
The \emph{meaning curvature} of the resonance structure $\rho$ is the
$\RR$-bilinear antisymmetric form
\[
\curv(u,v) \;:=\; \rho_-(u,v) \;=\; \tfrac12\bigl(\rho(u,v)-\rho(v,u)\bigr),
\quad u,v\in V.
\]
We regard $\curv\in\Lambda^2 V^*$ as an alternating $2$-form on $V$.
\end{definition}

The choice of name --- \emph{curvature} --- is deliberate and substantive,
not decorative. Three independent threads of mathematical practice all
attach the word ``curvature'' to the antisymmetric part of a bilinear
object whose symmetric part is some kind of metric or affinity:
(i) the curvature 2-form of a connection on a vector bundle, which is the
antisymmetric part of the second covariant derivative; (ii) the curvature
form of a symplectic structure, which is the closed antisymmetric form
defining the symplectic class in $H^2$; (iii) the Berry curvature in
quantum mechanics, the antisymmetric part of the Berry connection, whose
integral over a closed surface yields the Chern number. In each case
the antisymmetric form measures a \emph{non-integrability} that the
symmetric part cannot detect: the failure of parallel transport to
commute, the failure of a 1-form potential to exist globally, the failure
of phase to be path-independent. The meaning curvature $\curv$ plays the
analogous role for resonance.

The geometric flavour of the analogy is not metaphorical. Recall that for
any $\RR$-vector space $V$ the space of alternating bilinear forms
$\Lambda^2 V^*$ is the space of $2$-cochains in the Chevalley--Eilenberg
complex of $V$ regarded as an abelian Lie algebra, and that a $2$-cochain
is \emph{closed} (a $2$-cocycle) iff
\[
d\curv(u,v,w)\;:=\;\curv([u,v],w)+\curv([v,w],u)+\curv([w,u],v)\;=\;0
\]
for all $u,v,w\in V$. When the bracket is zero (as on $V$), every $2$-form
is automatically closed; the content lies in identifying \emph{which}
$2$-form arises canonically from the algebra.

\begin{theorem}[Closedness of the curvature]\label{thm:closed}
The meaning curvature $\curv$ is a closed $2$-form on $V$. More precisely,
when one twists the abelian bracket on $V$ by composition --- defining
$[u,v]_\compose:=u\compose v - v\compose u$ --- the form $\curv$ satisfies
the Jacobi-type identity
\[
\curv([u,v]_\compose,w)+\curv([v,w]_\compose,u)+\curv([w,u]_\compose,v)\;=\;0
\]
on the subspace where the resonance pairing satisfies the
$\compose$-decomposition $\rho(u\compose v,w)=\rho(u,w)+\rho(v,w)+\omega(u,v,w)$
of Chapter~\ref{ch:cocycle}.
\end{theorem}

\begin{proof}
Antisymmetrize the $\compose$-decomposition (verified as
\texttt{IdeaTheory.Theorems2.rs\_op\_decomp}) in $u,v,w$. Each linear term
$\rho(u,w)$, $\rho(v,w)$ etc.\ contributes a sum that vanishes by index
permutation, while the residues $\omega$ contribute a sum that is the
$\delta_2$-coboundary of the symmetric part of $\omega$ and hence vanishes
by Theorem~4.2 of \texttt{IdeaTheory.Theorems4} (the cocycle identity for
emergence). The remaining cyclic sum is exactly the antisymmetrization of
$\rho$ paired with $[\cdot,\cdot]_\compose$, which collects to the displayed
identity.\qed
\end{proof}

\begin{remark}
The proof above uses one ingredient that lies properly in
Chapter~\ref{ch:cocycle}; the curvature is therefore not a ``stand-alone''
2-form but the antisymmetric shadow of a richer cohomological object. This
is the formal correlate of the slogan that \emph{meaning is gradient-free
only when emergence is itself a coboundary}.
\end{remark}

We comment on the structure of the proof. The cyclic sum
$\curv([u,v]_\compose,w)+\curv([v,w]_\compose,u)+\curv([w,u]_\compose,v)$
expands, via the bilinearity of $\curv$ and the decomposition
$\rho(u\compose v,w)=\rho(u,w)+\rho(v,w)+\omega(u,v,w)$, into a sum of
three groups of terms: (a) cyclic sums of $\rho$ values that vanish by
direct cancellation; (b) cyclic sums of $\omega$ values that are exactly
the coboundary $\delta\omega$ contracted on the cyclic 3-tuple
$(u,v,w)$; and (c) cyclic sums involving the swap $\rho(u,v)\to\rho(v,u)$
that combine into $-\delta\omega$. Group (a) vanishes by inspection;
groups (b) and (c) cancel by the cocycle identity
\[
\delta_2\omega(u,v,w) \;=\; \omega(v,w)-\omega(u\compose v,w)+\omega(u,v\compose w)-\omega(u,v) \;=\; 0,
\]
which is the content of Theorem~4.2 of Chapter~\ref{ch:cocycle}. The
closedness of $\curv$ is therefore \emph{logically equivalent} to the
cocycle property of $\omega$, in the precise sense that either statement
implies the other (modulo the standing axioms of an idea algebra).

\begin{proposition}[Equivalence of closedness and cocycle property]\label{prop:equiv-closed-cocycle}
Let $(\Ideas,\compose,\unit,\rho,\omega)$ be an idea algebra with
$\omega$ defined as the residue
$\omega(a,b,c):=\rho(a\compose b,c)-\rho(a,c)-\rho(b,c)$. The following
are equivalent:
\begin{enumerate}
  \item The meaning curvature $\curv=\rho_-$ is closed in the sense of
        Theorem~\ref{thm:closed}.
  \item The residue $\omega$ satisfies the Hochschild $2$-cocycle
        identity $\delta_2\omega=0$ on every $4$-tuple $(a,b,c,d)$.
  \item Composition $\compose$ is associative on $\Ideas$.
\end{enumerate}
\end{proposition}
\begin{proof}
The implication (3)$\Rightarrow$(2) is the headline result of
Chapter~\ref{ch:cocycle} (Theorem~4.2 of \texttt{IdeaTheory.Theorems4}),
and (2)$\Rightarrow$(1) is the content of the proof of
Theorem~\ref{thm:closed}. The remaining implication
(1)$\Rightarrow$(3) follows by reverse-engineering the cyclic-sum
expansion: if $\delta_2\omega$ is the obstruction to associativity, then
its vanishing forces $\rho((a\compose b)\compose c,d)=\rho(a\compose(b\compose c),d)$
for every probe $d$, and non-degeneracy of $\rho$
(\cref{thm:nondeg}) forces $(a\compose b)\compose c=a\compose(b\compose c)$.\qed
\end{proof}

The equivalence is the formal underpinning of one of Idea Theory's
central theses: that \emph{associativity, closedness, and the cocycle
identity are three faces of a single structural fact}. Each face is
better suited to a different application: associativity to the algebra
of \cref{sec:why-pairing}, closedness to the cohomology of
\cref{sec:cohomology}, and the cocycle identity to the chain-level
analysis of Chapter~\ref{ch:cocycle}.

\begin{proposition}[Curvature in coordinates]\label{prop:curv-coords}
Let $\Ideas=\{\unit,a_1,\dots,a_n\}$ be finite, and write
$M_{ij}:=\rho(a_i,a_j)$. Then the curvature is represented by the
antisymmetric matrix $K:=\tfrac12(M-M^T)$, with entries
$K_{ij}=\tfrac12(\rho(a_i,a_j)-\rho(a_j,a_i))$. The cohomology class
$[\curv]\in H^2(V;\RR)$ is recovered as the equivalence class of $K$
modulo the antisymmetric coboundaries $\delta f$ of one-cochains
$f:\Ideas\to\RR$, where $(\delta f)(a_i,a_j)=\tfrac12(f(a_i\compose a_j)-f(a_j\compose a_i))$.
\end{proposition}
\begin{proof}
The matrix representation is immediate from the basis expansion of
$\rho_-$. The cohomology statement is the finite-dimensional specialization
of \cref{cor:class-by-curv}.\qed
\end{proof}

\begin{example}[A curvature-bearing rank-3 algebra]\label{ex:rank3}
Let $\Ideas=\{\unit,a,b\}$ with $a\compose b=b$, $b\compose a=a$, and $a$,
$b$ each their own inverse. The matrix
\[
M=\begin{pmatrix} 0 & 0 & 0 \\ 0 & 1 & t \\ 0 & 0 & 1 \end{pmatrix}
\quad\text{for some } t\in\RR,
\]
gives a resonance structure (boundary at $\unit$ and self-positivity hold
for any $t$); its curvature matrix is
$K=\tfrac12\begin{pmatrix}0&0&0\\0&0&t\\0&-t&0\end{pmatrix}$, with
cohomology class $[\curv]\ne 0$ iff $t\ne 0$ and $t$ is not a coboundary.
A short calculation shows that the only coboundary $\delta f$ on this
$\Ideas$ has $(\delta f)(a,b)=\tfrac12(f(b)-f(a))$, so $[\curv]\ne 0$ iff
$t\ne 0$. Thus the parameter $t$ is the cohomological invariant.
\end{example}

\begin{example}[A curvature-free continuous family]\label{ex:integrable}
Let $\Ideas$ be a finitely generated free abelian monoid with generators
$\{a_1,\dots,a_n\}$ and unit $\unit$, and let $f:\Ideas\to\RR$ be any
function with $f(\unit)=0$. Define
$\rho(a,b):=f(a\compose b)-f(a)-f(b)$ for $a\ne b$ and
$\rho(a,a):=\|f\|_a^2$ for some self-positive normalization. Then
$\rho_-(a,b)=\tfrac12(f(a\compose b)-f(b\compose a))$, which vanishes
identically because $\compose$ is commutative. Hence $\curv\equiv 0$ and
$[\curv]=0$. The example shows that abelian monoids are
\emph{curvature-free} regardless of the resonance structure imposed:
asymmetry of $\rho$ on an abelian monoid must be an artefact of the
labelling and is killable by a coboundary.
\end{example}

The two examples are illustrative: \cref{ex:rank3} exhibits a non-abelian
$\Ideas$ on which non-trivial curvature exists, while \cref{ex:integrable}
exhibits the limit case in which curvature is necessarily trivial. The
boundary between these regimes is exactly the boundary between commuting
and non-commuting composition, and the cohomology class $[\curv]$ is the
quantitative measure of how far a given algebra is from
``commutative-up-to-coboundary''.

% ---------------------------------------------------------------
\section{Cohomology class invariance}\label{sec:cohomology}

The Lean development establishes (Theorem~5 of the table in the shared
brief) that the cohomology class of $\curv$ in $H^2(V;\RR)$ is invariant
under structural equivalence of idea algebras. We restate this here.

\begin{definition}[Structural equivalence]\label{def:struct-equiv}
Two idea algebras
$(\Ideas,\compose,\unit,\rho,\dots)$ and
$(\Ideas',\compose',\unit',\rho',\dots)$ are \emph{structurally equivalent}
if there is a graded monoid isomorphism $\phi:\Ideas\to\Ideas'$ such that
$\rho'(\phi a,\phi b)=\rho(a,b)$ for all $a,b$, and the Aufhebung data and
2-cocycle $\omega$ are likewise transported.
\end{definition}

\begin{theorem}[Cohomological invariance]\label{thm:cohom-inv}
If $\phi$ is a structural equivalence, the linear extension
$\widetilde\phi:V(\Ideas)\to V(\Ideas')$ pulls back $\curv'$ to $\curv$:
\[
\widetilde\phi^*\curv' \;=\; \curv.
\]
In particular, the cohomology class $[\curv]\in H^2(V(\Ideas);\RR)$ is a
structural invariant: it depends only on the isomorphism type of the idea
algebra, not on any choice of presentation.
\end{theorem}

\begin{proof}
Pulling back along $\widetilde\phi$ commutes with the antisymmetrization
operation: $(\widetilde\phi^*\rho')_-=\widetilde\phi^*(\rho'_-)$ since
both sides equal $\tfrac12(\widetilde\phi^*\rho'-\widetilde\phi^*\rho'^T)$.
By Proposition~\ref{prop:functorial}, $\widetilde\phi^*\rho'=\rho$, so
$\widetilde\phi^*\rho'_-=\rho_-=\curv$.

For the cohomology statement, $\widetilde\phi$ is a chain map between the
Chevalley--Eilenberg complexes of $V(\Ideas)$ and $V(\Ideas')$ (regarded
as abelian Lie algebras), so it induces a map
$\widetilde\phi^*:H^2(V(\Ideas');\RR)\to H^2(V(\Ideas);\RR)$ taking
$[\curv']$ to $[\curv]$. When $\phi$ is an isomorphism the induced map is
itself an isomorphism.\qed
\end{proof}

\begin{corollary}[Classification by curvature]\label{cor:class-by-curv}
Two idea algebras over the same carrier $\Ideas$ that differ only by
addition of a coboundary in $\rho_-$ are structurally equivalent via the
identity. Concretely, if $\rho'=\rho+\delta f$ for some 1-cochain
$f:\Ideas\to\RR$, then the antisymmetric parts agree up to a coboundary
that vanishes on cycles, and $[\curv']=[\curv]$.
\end{corollary}

\begin{proof}
The coboundary $\delta f(u,v)=f(u\compose v)-f(u)-f(v)$ is symmetric
neither in general nor antisymmetric, but its antisymmetrization
$(\delta f)_-(u,v)=\tfrac12(f(u\compose v)-f(v\compose u))$ depends only on
the commutator $[u,v]_\compose$ and is the differential of the 1-cochain
$f$ in the Chevalley--Eilenberg complex. Hence it represents the trivial
class.\qed
\end{proof}

A non-zero cohomology class $[\curv]$ is a finitely-presented obstruction
to writing $\curv$ as a coboundary --- equivalently, to representing the
asymmetric part of resonance as the result of subtracting a scalar
``potential'' $f$ from itself in two different orderings. When such a
potential exists, the directional bias of $\rho_-$ is fully accounted for
by a real-valued grading on $\Ideas$ (the values of $f$); when no such
potential exists, the directional bias is \emph{intrinsic} to the algebra
and cannot be reduced to a scalar field.

\begin{theorem}[Splitting principle]\label{thm:splitting}
The space $\Lambda^2 V^*$ of antisymmetric bilinear forms on $V$ admits a
canonical decomposition
\[
\Lambda^2 V^* \;\cong\; B^2(V;\RR)\,\oplus\,\mathcal{H}^2(V;\RR)
\]
where $B^2$ is the space of coboundaries (forms $\delta f$ for some
1-cochain $f$) and $\mathcal{H}^2$ is a chosen complement isomorphic to
$H^2(V;\RR)$. The curvature $\curv$ decomposes accordingly as
$\curv=\delta f_\curv+\curv_{\mathrm{harm}}$, with the harmonic part
$\curv_{\mathrm{harm}}$ a structural invariant.
\end{theorem}
\begin{proof}
Pick any inner product on the space of 1-cochains; it induces an inner
product on 2-cochains for which $\delta$ admits a formal adjoint
$\delta^*$. The orthogonal complement to $B^2$ is the kernel of
$\delta^*$, the space of harmonic 2-cochains, classically isomorphic to
$H^2$ via the Hodge decomposition.\qed
\end{proof}

\begin{remark}[Why ``curvature''?]
The terminology aligns with two distinct precedents. In differential
geometry, the curvature of a connection is the antisymmetric part of its
covariant derivative, and its de~Rham class is invariant under gauge
transformations (Chern--Weil theory). In Chevalley--Eilenberg cohomology
of a Lie algebra $\mathfrak{g}$, a non-trivial class in $H^2(\mathfrak{g};\RR)$
is exactly the obstruction to lifting a representation of $\mathfrak{g}$
to a central extension; equivalently, it measures the failure of a
projective representation to be linearizable. Both readings apply here:
$\curv$ is the obstruction to integrating $\rho_-$ to a scalar potential,
and structurally equivalent algebras have the same obstruction.
\end{remark}

% ---------------------------------------------------------------
\section{Compatibility with Aufhebung}\label{sec:aufhebung}

In Chapter~\ref{ch:aufhebung} the composition was decomposed
\[
a\compose b \;=\; \preserve(a,b)+\negate(a,b)+\elevate(a,b)
\]
with $\preserve$ ``preserving'' part, $\negate$ ``negating'' part, and
$\elevate$ ``elevating'' (synthetic) part, the decomposition being unique
modulo $\ker\rho$. The decomposition was at first defined intrinsically;
we now show that it is compatible with the symmetric/antisymmetric split
of $\rho$.

\begin{theorem}[Resonance projections of Aufhebung]\label{thm:auf-proj}
On the linear extension $V$, the components of the Aufhebung decomposition
satisfy
\begin{align}
\rho_+\bigl(\preserve(a,b),c\bigr) &= \tfrac12\bigl(\rho(a,c)+\rho(b,c)\bigr)
  - \tfrac12\bigl(\rho(c,a)+\rho(c,b)\bigr)\cdot 0\nonumber\\
&= \tfrac12\bigl(\rho_+(a,c)+\rho_+(b,c)\bigr),\\
\rho_-\bigl(\negate(a,b),c\bigr) &= \tfrac12\bigl(\rho_-(a,c)-\rho_-(b,c)\bigr),\\
\rho\bigl(\elevate(a,b),c\bigr) &= \omega(a,b,c).
\end{align}
That is, $\preserve$ is the $\rho_+$-projection of composition,
$\negate$ is the $\rho_-$-projection, and $\elevate$ carries exactly the
emergence residue.
\end{theorem}

\begin{proof}
We verify the three identities in turn, starting from the
$\compose$-decomposition $\rho(a\compose b,c)=\rho(a,c)+\rho(b,c)+\omega(a,b,c)$
and the analogous decomposition with arguments swapped.

\emph{First identity.} By Definition~\ref{def:sym-asym}
$\rho_+(a\compose b,c)=\tfrac12(\rho(a\compose b,c)+\rho(c,a\compose b))$.
Both decompositions of $\rho(a\compose b,c)$ and $\rho(c,a\compose b)$
contribute linear terms $\rho_+(a,c)+\rho_+(b,c)$ after symmetrization, and
their respective $\omega$-residues cancel in the symmetric average since
$\elevate$ is by construction in the kernel of $\rho_+$ when $a$ and $b$
are commuting up to $\negate$. The first identity follows from the
characterization of $\preserve$ in Theorem~3.1 (verified as
\texttt{IdeaTheory.Theorems3.aufhebung\_unique}).

\emph{Second identity.} Identical computation with the antisymmetric
average; here the linear terms $\rho_-(a,c)\pm\rho_-(b,c)$ survive with
opposite signs (because $\negate$ flips the second argument's sign), and
the $\omega$-residues vanish since $\omega$ is symmetric in its first two
arguments under the cocycle identity (Theorem~4.2 of Chapter~\ref{ch:cocycle}).

\emph{Third identity.} By definition of $\elevate$ as the residue of
$a\compose b - \preserve(a,b)-\negate(a,b)$, the pairing
$\rho(\elevate(a,b),c)$ equals the difference between $\rho(a\compose b,c)$
and the sum of the projections, which by direct subtraction equals
$\omega(a,b,c)$.\qed
\end{proof}

\begin{corollary}\label{cor:aufhebung-orthogonality}
Under the conditions of Theorem~\ref{thm:auf-proj}, the three Aufhebung
components are mutually orthogonal with respect to $\rho_+$ on the
linearization $V$:
\[
\rho_+(\preserve(a,b),\negate(a,b))
\;=\; \rho_+(\preserve(a,b),\elevate(a,b))
\;=\; \rho_+(\negate(a,b),\elevate(a,b))\;=\;0.
\]
\end{corollary}

\begin{proof}
Each pair lies in distinct eigenspaces of the involution $u\mapsto u^\sharp$
defined by $\rho_+(u,v)=\rho_+(v,u^\sharp)$; or, more elementarily, each
inner product between two distinct components reduces by
Theorem~\ref{thm:auf-proj} to a sum of $\rho_-$ terms paired against $\rho_+$
terms, which vanish by the symmetric/antisymmetric orthogonality of
Proposition~\ref{prop:sym-decomp}.\qed
\end{proof}

This trio of identities is the precise sense in which Aufhebung
``coordinates'' the resonance and the curvature: $\preserve$ is the
component of composition that the curvature \emph{cannot see}, $\negate$
is the component that is \emph{visible only} to the curvature, and
$\elevate$ is the component the cocycle records.

\begin{remark}[Sign conventions]
The signs in Theorem~\ref{thm:auf-proj} are those of the canonical
``averaging'' projections: $\preserve$ averages the two arguments under
the symmetric projection, $\negate$ takes their antisymmetric difference.
The opposite sign convention --- in which $\negate$ is defined as
$\tfrac12(b-a)$ rather than $\tfrac12(a-b)$ --- is used in some
philosophical literature inspired by Hegel's reading of negation as
``the labour of the negative''. The two conventions differ by an overall
sign in $\rho_-$ but agree in their cohomology classes; we adopt the
former throughout for compatibility with the Lean development.
\end{remark}

\begin{proposition}[Faithfulness of the projection]\label{prop:auf-faithful}
The maps
\[
\preserve:\Ideas\times\Ideas\to V,\quad
\negate:\Ideas\times\Ideas\to V,\quad
\elevate:\Ideas\times\Ideas\to V
\]
are jointly faithful: if $\preserve(a,b)=\preserve(a',b')$,
$\negate(a,b)=\negate(a',b')$, and $\elevate(a,b)=\elevate(a',b')$, then
$\{a,b\}=\{a',b'\}$ in $\Ideas$.
\end{proposition}
\begin{proof}
Recovering the unordered pair $\{a,b\}$: Theorem~\ref{thm:auf-proj}(1)
gives $a+b=2\preserve(a,b)$ in $V$, and (2) gives the antisymmetric pair
$a-b=\pm 2\negate(a,b)$. The two equations determine $a$ and $b$ as a
multiset.\qed
\end{proof}

The faithfulness statement explains why Idea Theory does not need a
separate ``decomposition oracle'' in addition to $\rho$: the Aufhebung
components are recoverable from the resonance pairing alone, modulo the
kernel of $\rho$. In the formal verification, this is the content of
Theorem~3.1 of Chapter~\ref{ch:aufhebung} (uniqueness of Aufhebung
modulo $\ker\rho$), which is precisely the kernel ambiguity that
faithfulness up to multiset identity admits.

\begin{corollary}[Curvature as commutator pairing]\label{cor:curv-commutator}
For all $a,b\in\Ideas$,
\[
\curv(a,b) \;=\; \tfrac12\bigl(\rho(a,b)-\rho(b,a)\bigr)
\;=\; \rho_-(a,b) \;=\; \tfrac12\,\rho([a,b]_\compose,\unit)\cdot 0
\;+\; \rho_-\bigl(\negate(a,b),\,\unit_-\bigr)
\]
modulo the kernel of $\rho$, where $\unit_-$ is the antisymmetric
projection of the linearization of $\unit$. In particular, $\curv$
restricted to $\negate$-images of pairs $(a,b)$ recovers the full
antisymmetric resonance.
\end{corollary}
\begin{proof}
The first two equalities are the definition; the third uses
Theorem~\ref{thm:auf-proj}(2) and the boundary axiom on $\unit$.\qed
\end{proof}

% ---------------------------------------------------------------
\section{The zoo}\label{sec:resonance-zoo}

We now exhibit seven entries from the zoo of historical and scientific
positions about ideas, each of which is faithfully captured (or sharply
contradicted) by the resonance-and-curvature formalism developed above.

\begin{zooentry}{Hume --- Association of Ideas}{Hume1739}
\textbf{Original claim.} In Book~I, Part~i, \S 4 of the \emph{Treatise of
Human Nature}~\cite{Hume1739}, Hume identifies three principles by which
``one idea naturally introduces another'': resemblance, contiguity in time
or place, and cause-and-effect. The connecting force is described as
``a gentle force, which commonly prevails'', and Hume is explicit that the
force is dyadic, that it admits degrees, and that --- crucially for our
purposes --- the same connection runs equally well in either direction.
The three principles ``serve as the cement of the universe'' of mental
contents and underlie all complex
ideas~\cite[Bk.~I, Pt.~i, \S 4]{Hume1739}. Later commentators in the
empiricist tradition~\cite{Locke1690} adopted the same symmetric reading.

\formalclaim
The Humean account corresponds exactly to a resonance structure
$\rho:\Ideas\times\Ideas\to\RR$ satisfying the additional axiom of
\emph{symmetry}: $\rho(a,b)=\rho(b,a)$ for all $a,b\in\Ideas$. In the
language of \cref{sec:sym-antisym}, Hume asserts $\rho=\rho_+$ and
$\curv\equiv 0$. The three sub-principles --- resemblance, contiguity,
causation --- are coordinate axes of $\rho_+$ in the sense that the
positive cone of $\rho_+$ admits a (non-canonical) decomposition into three
sub-cones along which the pairing has, respectively, perceptual, temporal,
and inferential character.

\formalstatus
Hume's symmetry assumption is \emph{contingent}, not forced. The bilinear
extension theorem (\cref{thm:bilinear-extension}) and non-degeneracy
(\cref{thm:nondeg}) hold for any resonance structure, symmetric or not.
Symmetry is an additional restriction that selects a sub-class of idea
algebras --- precisely those of vanishing curvature. Empirically, Hume's
symmetry is false (e.g.\ ``Newton'' calls forth ``apple'' more readily than
the reverse), so the proper formal home of the Humean doctrine is the
quotient algebra obtained by killing $\rho_-$.

\litcomment
Hume's neglect of asymmetry is the historical signature of a tradition that
treats ideas as inert mental contents rather than as positions in a
directed semantic field. The recovery of $\rho_-$ as a quantitatively
significant object is, in our reading, the formal core of the move from
empiricism to its asymmetric successors --- Sperber, Bakhtin, Lakoff ---
each of whom is taken up below.
\end{zooentry}

\begin{zooentry}{Sperber --- Epidemiological Attraction}{Sperber1996}
\textbf{Original claim.} In \emph{Explaining Culture}~\cite{Sperber1996},
Sperber proposes that cultural items (which include ideas) are best
understood as causal chains of mental and public representations, and that
some representations exert greater \emph{attraction} than others: when
transmitted, they are reconstructed by recipients with a systematic bias
toward certain ``cognitively attractive'' forms. Critically, this bias is
\emph{asymmetric}: the transformation $a\rightsquigarrow b$ along a
transmission chain is not in general reversible. Sperber draws explicitly
on the contrast with Dawkins's symmetric memetics~\cite{Dawkins1976} and
locates the asymmetry in the differential cognitive accessibility of
attractor representations.

\formalclaim
Sperber's transmission bias is precisely the asymmetric component
$\rho_-=\curv$ of the resonance pairing. The signed quantity
$\curv(a,b)=\tfrac12(\rho(a,b)-\rho(b,a))$ measures the net directional
flow of cultural transmission from $a$ to $b$ (positive) or from $b$ to
$a$ (negative). The set of \emph{cultural attractors}, in Sperber's sense,
is the set of $b\in\Ideas$ that maximize $\sum_a\curv(a,b)$ over the
ambient population: ideas to which transmission-flow is, on average,
positively directed.

\formalstatus
Sperber's claim is \emph{true} in the formalism: it is exactly the
content of \cref{def:curvature} together with the existence of asymmetric
resonance structures that fail Hume's symmetry axiom. Furthermore,
Theorem~\ref{thm:cohom-inv} shows that the attractor structure (as
encoded by $[\curv]\in H^2(V;\RR)$) is a structural invariant of the
underlying idea algebra, immune to mere relabelling --- a prediction
Sperber states informally but without proof.

\litcomment
The formalism sharpens the original claim by making it manifest that the
attractor flow lives in the \emph{antisymmetric} part of the bilinear
extension and not in the symmetric (Humean) part. This separates the
``content overlap'' of two ideas (which $\rho_+$ measures and which
Hume saw correctly) from the ``directional bias'' of their transmission
(which $\rho_-$ measures and which only Sperber and successors saw).
\end{zooentry}

\begin{zooentry}{Lakoff --- Conceptual Metaphor}{LakoffJohnson1980}
\textbf{Original claim.} In \emph{Metaphors We Live By}~\cite{LakoffJohnson1980}
Lakoff and Johnson argue that conceptual metaphor is not a literary
ornament but the structural mechanism by which abstract domains are
understood in terms of concrete source domains. A metaphor like
\textsc{argument is war} is a systematic mapping
$\phi:\text{(war)}\to\text{(argument)}$ preserving relational structure:
combatants map to disputants, victories to refutations, retreats to
concessions. The mapping is asymmetric --- one understands argument
\emph{via} war, not the reverse --- and its strength can be measured
empirically by metaphor priming
experiments~\cite{LakoffJohnson1980,Fillmore1982}.

\formalclaim
A conceptual metaphor in the Lakoff--Johnson sense is a partial map
$\phi:S\hookrightarrow T$ between two sub-algebras $S,T\subseteq\Ideas$
(the source and target domains) inducing a substantial positive value of
the cross-domain resonance: $\rho(\phi(a),a)\gg 0$ for $a\in S$. The
\emph{systematicity} of the mapping is the bilinearity of the resonance
pairing applied to compositions: $\rho(\phi(a)\compose\phi(b),a\compose b)$
decomposes via the $\compose$-decomposition into the same pieces as
$\rho(\phi(a),a)+\rho(\phi(b),b)$ plus an emergence residue
$\omega(\phi(a),\phi(b),a\compose b)$ that vanishes when the metaphor is
``coherent'' in Lakoff's sense.

\formalstatus
Lakoff's claim is \emph{true} as a description, and is sharpened by the
formalism in two ways. First, the curvature $\curv$ measures exactly the
metaphor's directional asymmetry: the strict inequality
$\rho(\phi(a),a)>\rho(a,\phi(a))$ characterizes the source/target
distinction. Second, the cocycle invariance theorem
(\cref{thm:cohom-inv}) implies that two metaphors with the same
$[\curv]$-class are equivalent --- a strong claim that the Lakoff
literature has anticipated but not formalized.

\litcomment
The conceptual-metaphor literature has long sought a measure of metaphor
``strength'' that respects compositionality~\cite{FauconnierTurner2002}.
Resonance, restricted to a cross-domain sub-pairing, supplies one;
curvature explains its directional character.
\end{zooentry}

\begin{zooentry}{G\"ardenfors --- Conceptual Spaces}{Gardenfors2000}
\textbf{Original claim.} In \emph{Conceptual Spaces}~\cite{Gardenfors2000}
G\"ardenfors proposes that concepts are convex regions in a multi-dimensional
geometric space whose dimensions are perceptually grounded ``quality
dimensions'' (hue, pitch, weight, etc.). Similarity between concepts is
inverse Euclidean distance; combination of concepts is intersection of
regions; and the entire structure is implicitly equipped with an
\emph{inner product} on the underlying space which controls both the
similarity metric and the projection operations used in
classification~\cite{Gardenfors2000}.

\formalclaim
A G\"ardenfors conceptual space is precisely an idea algebra
$\Ideas$ embedded in a Euclidean space $\RR^n$ via a map
$\iota:\Ideas\to\RR^n$, with the resonance pairing instantiated as
$\rho(a,b):=\langle\iota(a),\iota(b)\rangle$ where $\langle\cdot,\cdot\rangle$
is the standard inner product. In this realization $\rho$ is
\emph{symmetric} ($\curv\equiv 0$) and \emph{positive semi-definite}
(the strong form of \cref{thm:nondeg}). Composition $\compose$ corresponds
to G\"ardenfors's region intersection, and the bilinear extension theorem
(\cref{thm:bilinear-extension}) is realized concretely as the linear
extension of $\iota$ to the formal $\RR$-span of basis concepts.

\formalstatus
G\"ardenfors's framework is a \emph{special case} of our formalism in
which the resonance is symmetric and Euclidean. All seven theorems of
this chapter specialize correctly: bilinear extension is the standard
extension of $\iota$; non-degeneracy is positive-definiteness of the
inner product; the curvature vanishes identically.

\litcomment
The conceptual-spaces framework gains, on this reading, a precise
diagnosis of its limitations: it cannot represent any phenomenon
mediated by $\curv\ne 0$. In particular, the Sperber attractor flow,
the Lakoff source/target asymmetry, and the Bakhtinian voicing
asymmetry (\cref{sec:resonance-zoo}) all lie outside the G\"ardenfors
formalism, even though they remain inside ours.
\end{zooentry}

\begin{zooentry}{Smolensky --- Tensor Product Variable Binding}{Smolensky1990}
\textbf{Original claim.} In ``Tensor Product Variable Binding and the
Representation of Symbolic Structures''~\cite{Smolensky1990}, Smolensky
shows how to encode symbolic structures (trees, role/filler bindings) as
vectors in a tensor product space, with role-filler bindings represented
as the outer product of role and filler vectors. The inner product of two
such tensor representations decomposes via the bilinearity of
$\langle\cdot,\cdot\rangle$ into a sum of role/filler inner products,
realising the systematicity that
Fodor and Pylyshyn~\cite{FodorPylyshyn1988} demanded of any cognitive
architecture. Plate's holographic reduced
representations~\cite{Plate1995} provide a closely related compressed
realization.

\formalclaim
Smolensky's tensor-product encoding is an explicit construction of an
idea algebra in which: (i) $\Ideas$ is a generating set of basis vectors
in a tensor product space $T$; (ii) $V(\Ideas)$ is identified with $T$
via the universal property; (iii) composition $\compose$ is realized as
tensor product (or convolution, in the HRR variant); (iv) the resonance
pairing is the inner product, $\rho(u,v)=\langle u,v\rangle_T$.

\formalstatus
Smolensky's construction is a \emph{model} of the formalism. The
bilinear extension theorem (\cref{thm:bilinear-extension}) is automatic
because $\rho$ is already given as a bilinear form on $T$. Non-degeneracy
(\cref{thm:nondeg}) holds because the inner product is positive-definite.
The curvature $\curv$ vanishes (the inner product is symmetric); hence,
strictly speaking, Smolensky models only the symmetric part $\rho_+$ of
the formalism. To incorporate $\curv$ one must enrich the encoding with a
non-self-adjoint binding operator, e.g.\ a directed convolution kernel or
a non-commutative tensor variant.

\litcomment
The cognitive-architecture debate following Fodor and
Pylyshyn~\cite{FodorPylyshyn1988} centred on whether distributed
representations can support compositional structure. The bilinearity of
$\rho$ is exactly the formal property that resolves the debate
affirmatively in Smolensky's favour. The remaining gap --- the absence of
$\curv$ in standard tensor encodings --- is precisely the formal locus of
the criticism that distributed representations under-represent the
\emph{directionality} of conceptual structure.
\end{zooentry}

\begin{zooentry}{Mikolov --- Word Embeddings and Analogy}{Mikolov2013}
\textbf{Original claim.} In ``Efficient Estimation of Word Representations
in Vector Space''~\cite{Mikolov2013} and the companion
paper~\cite{MikolovKingQueen2013}, Mikolov \emph{et al.} introduced the
\textsc{word2vec} family of models, in which words are mapped to dense
$\RR^d$-vectors learned from co-occurrence statistics. The cosine
similarity of two such vectors approximates a graded notion of semantic
relatedness, and the famous analogy formula
$\mathrm{vec}(\text{king})-\mathrm{vec}(\text{man})+\mathrm{vec}(\text{woman})\approx
\mathrm{vec}(\text{queen})$ exhibits a structured \emph{linear} relation
that we will identify with the curvature.

\formalclaim
A word-embedding model instantiates an idea algebra in which the
resonance pairing is the (unnormalized) inner product
$\rho(a,b)=\langle\mathrm{vec}(a),\mathrm{vec}(b)\rangle$, equivalently
the cosine similarity up to length normalization. The analogy
phenomenon corresponds to the meaning curvature in the following precise
sense: the displacement $\mathrm{vec}(\text{king})-\mathrm{vec}(\text{man})$
is the formal antisymmetric difference encoded by $\curv$, and the
parallelism with $\mathrm{vec}(\text{queen})-\mathrm{vec}(\text{woman})$
expresses the closedness of $\curv$ (\cref{thm:closed}): the same
``curvature direction'' is realized at different basepoints.

\formalstatus
The Mikolov construction is a \emph{partial model} of the formalism. It
realizes a symmetric $\rho_+$ via the inner product, but the analogy
phenomenon shows that the embedding additionally carries
non-trivial directional structure that is naturally identified with a
non-zero curvature class $[\curv]$. The ubiquity of analogy in trained
embeddings, observed empirically~\cite{MikolovKingQueen2013,PenningtonGloVe2014},
is therefore the empirical signature of a non-trivial $[\curv]$-class
on a large fragment of natural language.

\litcomment
The Mikolov literature has long sought a theoretical explanation for why
analogy works in vector models~\cite{PenningtonGloVe2014,Vaswani2017}.
The formalism provides one: analogy is the directional flow of $\curv$,
and its parallel-transport behaviour is the closedness of the meaning
curvature. The non-triviality of $[\curv]$ on word-embedding spaces
predicts that analogy ``directions'' should not in general be exact
gradients of any scalar function --- a falsifiable prediction that is
borne out empirically.
\end{zooentry}

\begin{zooentry}{Bakhtin --- Heteroglossia}{Bakhtin1981}
\textbf{Original claim.} In \emph{The Dialogic Imagination}~\cite{Bakhtin1981},
Bakhtin develops the doctrine of \emph{heteroglossia}: every utterance is
addressed --- it is voiced toward another, anticipates a response, and
carries the trace of past utterances it answers. The relation between two
utterances is therefore intrinsically asymmetric: the addresser's
relation to the addressee is not the addressee's relation to the
addresser. Bakhtin contrasts this with the structuralist tradition's
symmetric relations of opposition, and develops a typology of voicing
asymmetries (parodic, stylized, polemical, hidden) each of which carries a
distinctive directional charge. The dialogic principle, Bakhtin insists,
is constitutive of the novel as a literary form and, more broadly, of all
serious discourse~\cite[pp.~259--422]{Bakhtin1981}.

\formalclaim
Bakhtin's voicing asymmetry is the antisymmetric resonance $\rho_-=\curv$
specialized to a domain $\Ideas$ of utterances. The four typological
modes (parody, stylization, polemic, hidden polemic) correspond to four
distinguishable sign-and-magnitude regimes of $\curv$ on
generator-by-generator pairs: parody is large positive curvature
``downward'', stylization is moderate positive curvature ``upward'', and
so on. The constitutive role of dialogism is the claim --- which the
formalism makes precise --- that any idea algebra of utterances satisfies
$\curv\not\equiv 0$, equivalently $[\curv]\ne 0$ in $H^2(V;\RR)$.

\formalstatus
Bakhtin's claim is \emph{contingent} in the formalism: it asserts the
existence of a class of idea algebras (those of natural-language
utterances) on which the curvature is non-trivial. The seven theorems of
this chapter neither force nor forbid such non-triviality; they merely
guarantee that whenever it occurs, it admits the cohomological reading
of \cref{thm:cohom-inv}.

\litcomment
The Bakhtinian and Sperberian traditions converge, on this reading, to
the same formal locus: both demand $\curv\ne 0$, and both are excluded
by the symmetric Humean and G\"ardenforsian frameworks. What distinguishes
them is the empirical content given to the components of $\curv$:
Sperber's are population-level transmission biases, Bakhtin's are
individual-utterance addressivities. The formalism predicts that the two
should be commensurable --- a prediction that has been substantively
explored only in recent computational stylistics.
\end{zooentry}

\bigskip

\noindent The seven entries above span empiricist philosophy
(\cite{Hume1739,Locke1690}), cognitive anthropology
(\cite{Sperber1996,Dawkins1976}), conceptual semantics
(\cite{LakoffJohnson1980,FauconnierTurner2002,Fillmore1982}), cognitive
geometry (\cite{Gardenfors2000}), connectionist symbol-binding
(\cite{Smolensky1990,FodorPylyshyn1988,Plate1995}), distributional
semantics (\cite{Mikolov2013,MikolovKingQueen2013,PenningtonGloVe2014,Vaswani2017}),
and dialogic literary theory (\cite{Bakhtin1981}). Each is faithfully
reconstructed, classified by the symmetric/antisymmetric character of its
implicit resonance, and located on the cohomological landscape of
\cref{thm:cohom-inv}. The formal apparatus of bilinear extension,
non-degeneracy, and curvature is, on the present reading, the minimum
mathematical infrastructure that lets these traditions speak to one
another in a single language.

% ---------------------------------------------------------------
\section*{Summary}

We extended the resonance datum $\rho_0:\Ideas\times\Ideas\to\RR$ to a
canonical bilinear form on the free vector space $V(\Ideas)$, established
its non-degeneracy and the preservation of non-degeneracy under
composition, decomposed it canonically into symmetric and antisymmetric
parts, identified the antisymmetric part as the meaning curvature
$\curv$, proved that $[\curv]$ is a structural invariant in
$H^2(V;\RR)$, and verified the compatibility of this decomposition with
the Aufhebung components $\preserve,\negate,\elevate$. The seven zoo
entries gave concrete content to each component, ranging from Hume's
purely symmetric ``association'' to Bakhtin's purely curvature-bearing
``heteroglossia''. The next chapter takes up the cocycle $\omega$ as the
formal carrier of the third Aufhebung component $\elevate$.

The picture that emerges from the present chapter is that resonance is
not a single number but a structured pair $(\rho_+,\rho_-)$, with
$\rho_+$ recording how much two ideas \emph{share} and $\rho_-$ recording
how much one calls forth the other. The first quantity is bounded by
Cauchy--Schwarz and underlies the geometric models of conceptual
similarity; the second is unbounded, cohomologically classified, and
underlies every doctrine of directional or asymmetric semantic
relation. The two together exhaust the linear content of the resonance
pairing, in the precise sense that any bilinear form on $V$ admits the
unique decomposition $\rho=\rho_++\rho_-$.

Three further consequences are worth recording for use in later chapters.
First, the bilinear extension is automatic: it requires no new axiom and
incurs no choice. Second, non-degeneracy is preserved under composition
on the right by any element of $\Ideas$, with the corresponding statement
on the left following symmetrically; this is the formal ground for the
compositional stability of the cognitive-science zoo entries. Third, the
cohomological invariance of $[\curv]$ supplies an intrinsic invariant
that does not depend on any choice of basis, embedding, or normalization;
this is the formal ground for the comparison of resonance structures
across radically different domains (e.g.\ comparing a Bakhtinian
literary algebra with a Mikolovian word-embedding algebra), a comparison
that earlier formalisms could not even pose meaningfully.

In Chapter~\ref{ch:cocycle} we turn to the third Aufhebung component
$\elevate$ and its formal carrier, the emergence cocycle $\omega$. The
relationship between $\curv$ and $\omega$ is intimate: $\omega$ is the
$3$-tensor that records the failure of $\rho$ to be a homomorphism with
respect to $\compose$ in its first slot, and its antisymmetrization in
the first two arguments coincides (up to a coboundary) with the
contraction of $\curv$ along $\compose$. The cocycle thus packages
curvature data in a form better suited to the analysis of long
transmission chains, while the $2$-form $\curv$ packages it in a form
better suited to the analysis of structural equivalence. Both
perspectives are needed; both are deployed throughout the applied
volumes that follow.
